\begin{abstract}
Ce rapport rassemble, en un seul document, l'état du projet « jumeau numérique de l'offre de soins libérale » au
15 novembre 2026, pour les équipes qui le reprendront : chercheurs, statisticiens et professionnels de la
planification sanitaire.

\paragraph{Le modèle.} On projette jusqu'en 2050 l'effectif du secteur privé (« libéral » au sens large) de cinq
professions : sages-femmes, chirurgiens-dentistes, pharmaciens, masseurs-kinésithérapeutes et infirmiers. Le
calcul descend à la commune, mais l'affichage et la validation s'arrêtent au département. Le modèle simule des
politiques publiques nationales exprimées comme des réglages datés d'hyperparamètres dont la source et le statut
(mesuré, sourcé ou supposé) sont déclarés. C'est un jumeau à état explicite (profession × territoire × âge × sexe ×
secteur), composé de modules de flux à contrat vérifiable. Les critères d'acceptation du jumeau et l'analyse de
sensibilité (axe~5) portent sur sa composition de référence, le stock-flux A0. L'estimateur servi par la carte et
l'assistant est un modèle de cohortes (C1), branché dans le jumeau par injection de ses taux de rétention et
d'entrée.

\paragraph{Le principal résultat de précision.} Au banc départemental à cinq ans (protocole N1, secteur privé,
\num{2056} couples département × cellule), l'erreur absolue moyenne en pourcentage de C1 est de
7,76\,\%\mesure{} [7,12 ; 8,40], contre 9,94\,\%\mesure{} pour la persistance et 9,32\,\%\mesure{} pour la
tendance linéaire. Ce banc est rétrospectif, et une partie de ses cibles a été vue à la conception (le pilote de
C1 a vu deux cellules de 2021) : le seul test propre est le dépôt de prévisions jugé en 2027. Aucune bande de C1
n'est calibrée (couverture de 73,5\,\%\mesure{} pour une cible de 75 à 85\,\%, protocole N6). Le budget d'erreur
attribue au bruit de comptage un plancher de 3,11 points\derive{} [2,89 ; 3,36], soit environ 40\,\%\derive{} de
la MAPE ; c'est une borne basse de la part irréductible. Trois familles de pistes n'ont pas battu C1 : des
intervalles de prévision recalibrés (prévision conforme, réconciliation MinT), un gradient boosting global et un
perceptron multicouche ; ces deux derniers sont rejetés contre A2 selon le test inscrit, et contre C1 à titre
exploratoire.

\paragraph{Ce qui n'est pas acquis.} Plusieurs résultats sont négatifs et sont exposés comme tels :
\begin{itemize}
\item aucun chiffre à l'horizon 2050 n'est citable comme prévision ;
\item aucun grain inférieur au département n'est calibré ;
\item aucun levier n'est validé de bout en bout par une évaluation causale ;
\item pour la moitié des protocoles pré-enregistrés, l'antériorité du gel sur le calcul n'est pas prouvable par
  l'historique du dépôt.
\end{itemize}

\paragraph{Le reste du document.} Il couvre aussi l'interface, la couche de langage naturel et l'assistant, la
conformité (RGPD, AIPD, règlement 2017/745), un inventaire chiffré des questions ouvertes et le programme du
module de besoin de soins. Chaque nombre porte sa provenance : mesuré\mesure, dérivé\derive{} ou
supposé\suppose. Les incohérences relevées entre les fichiers du dépôt sont listées en annexe, avec la valeur
retenue.
\end{abstract}
